\chapter{Цифрова модуляція: практичне ядро для SDR/DSP}
\label{ch:digital-modulation-practical-core}

\sourcebox{Основа: вже наявна українська доріжка PySDR \texttt{content-ukraine}, з нормалізацією у TeX і додаванням формульної рамки. Це не field-configuration recipe; це навчальний модуль про символи, сузір'я, I/Q-представлення і базову смугу.}

\section{Символи, біти і спектральна ефективність}

Цифрова модуляція передає інформацію не безпосередньо як ``текст'', а як послідовність \emph{символів}. Символ - це один елемент скінченного набору сигналів. Якщо сузір'я має $M$ можливих символів, то один символ несе
\[
  b = \log_2 M
\]
бітів, якщо $M$ є степенем двійки і всі символи використовуються рівноймовірно. Якщо символи передаються з швидкістю $R_s$ symbols/s, то ідеальна бітова швидкість до кодування дорівнює
\[
  R_b = R_s \log_2 M.
\]
Практичний компроміс: збільшення $M$ підвищує spectral efficiency, але зменшує відстань між точками сузір'я при фіксованій потужності. Тому сигнал стає чутливішим до шуму, phase error, frequency offset і нелінійностей підсилювача.

\section{Baseband та passband}

У SDR найзручніше мислити у комплексній базовій смузі. Нехай комплексний baseband-сигнал
\[
  x(t) = I(t) + j Q(t).
\]
Реальний passband-сигнал з несучою частотою $f_c$ можна записати як
\[
  s(t)=\Re\{x(t)e^{j2\pi f_c t}\}
      = I(t)\cos(2\pi f_c t)-Q(t)\sin(2\pi f_c t).
\]
Ця формула є головним мостом між математикою DSP і фізичним RF-ланцюгом: I/Q-дані живуть у цифровому baseband, а антена випромінює реальний passband-сигнал.

\section{ASK, PSK, FSK, QAM}

\subsection{ASK}
Amplitude-shift keying змінює амплітуду символу. Для $M$-ASK множина символів може бути
\[
  \mathcal S_{ASK}=\{a_1,a_2,\ldots,a_M\}\subset\R.
\]
Перевага - простота. Недолік - висока чутливість до amplitude fading і gain error.

\subsection{PSK}
Phase-shift keying тримає амплітуду сталою і змінює фазу:
\[
  s_k = e^{j2\pi k/M},\qquad k=0,1,\ldots,M-1.
\]
Для BPSK маємо дві точки $\{+1,-1\}$. Для QPSK - чотири фази; типово один символ несе 2 bits.

\subsection{QAM}
Quadrature amplitude modulation змінює і $I$, і $Q$. Для квадратного $M$-QAM точки лежать на ґратці в комплексній площині:
\[
  s_{m,n}=a_m + j a_n.
\]
QAM ефективний за спектром, але потребує кращого SNR і акуратнішої синхронізації. Для normalized constellation середню енергію символу часто приводять до
\[
  E_s = \mathbb E[|s|^2]=1.
\]

\subsection{FSK}
Frequency-shift keying кодує символи різними частотами. Для ортогонального FSK різницю частот обирають так, щоб інтеграл взаємного добутку символів за інтервал символу був нульовий. На практиці FSK може бути робастним, але зазвичай менш spectral-efficient.

\section{Сузір'я та мінімальна відстань}

Constellation diagram показує множину можливих $I/Q$-символів. Для AWGN-каналу груба помилка часто визначається мінімальною евклідовою відстанню
\[
  d_{\min}=\min_{i\ne k}|s_i-s_k|.
\]
Більший $d_{\min}$ при тій самій енергії означає меншу ймовірність помилки. Тому modulation design постійно балансує між $\log_2 M$, bandwidth, transmit power, robustness і hardware limits.

\section{Імпульсне формування}

Символи не можна просто передавати як прямокутні імпульси без спектральної ціни: різкі фронти мають широкий спектр. Тому вводять shaping pulse $p(t)$:
\[
  x(t)=\sum_{k} a_k p(t-kT_s),
\]
де $a_k$ - комплексні символи, $T_s$ - symbol period. Класична умова відсутності міжсимвольної інтерференції у моменти sampling:
\[
  p(nT_s)=\begin{cases}
  1, & n=0,\\
  0, & n\ne 0.
  \end{cases}
\]
Raised-cosine і root-raised-cosine filters використовуються саме для керування смугою та ISI.

\section{Практичні перевірки для перекладача}

\begin{enumerate}[label=\arabic*.]
\item Не перекладати \sourcepath{I/Q} як ``я/питання''. Це in-phase/quadrature, тобто синфазна і квадратурна компоненти.
\item \textit{Symbol} у цьому контексті - \emph{символ}, не ``знак'' у типографському сенсі.
\item \textit{Constellation} - \emph{сузір'я}; перше вживання бажано пояснити як діаграму можливих I/Q-точок.
\item \textit{Baseband} у цьому корпусі фіксуємо як \emph{базова смуга}; \textit{passband} - \emph{смуга перенесення} або \emph{смуговий сигнал} залежно від речення.
\item Формули з $2\pi f_c t$ не міняти на $\omega_c t$, якщо source цього не робить; якщо source має $\omega_c$, тоді $\omega_c=2\pi f_c$ треба явно вказати.
\end{enumerate}

